\section{Algorithm}
\label{sec:algorithm}

\subsection{ShortBurstForest}
\label{subsec:alg-sbf}

\sbf{} is the \sburst{} meta-algorithm with ForestRecom as the inner chain.
Algorithm~\ref{alg:sbf} gives the complete specification.

\begin{algorithm}
\caption{\sbf{}}
\label{alg:sbf}
\begin{algorithmic}[1]
\Require Burst length $\ell$, burst count $n$, percentile $p \in [0,1]$,
         base seed $b$
\Ensure  Selected plan $\pi^*$
\State $\pi_{\mathrm{cur}} \gets \mathrm{METIS\_plan}(b)$
       \Comment{initial plan from METIS}
\State $\mathrm{endpoints} \gets []$
\For{$i \gets 0$ \textbf{to} $n-1$}
  \State $s_i \gets \mathrm{SHA256}\!\left(\texttt{"SBF\_CHAIN\_"} \,\|\,
         i^{\mathrm{le64}} \,\|\, \texttt{"\_"} \,\|\, b^{\mathrm{le64}}
         \right)[0{:}8]\;\text{as } u64$
  \State $\mathrm{chain} \gets \mathrm{ForestRecomChain}(\pi_{\mathrm{cur}}, s_i)$
  \State $\mathrm{accepted}_i \gets 0$
  \For{$t \gets 0$ \textbf{to} $\ell - 1$}
    \State $s_{\mathrm{fwd}} \gets \mathrm{SHA256}\!\left(\texttt{"SBF\_FWD\_"} \,\|\,
           t^{\mathrm{le32}} \,\|\, \texttt{"\_"} \,\|\, s_i^{\mathrm{le64}}
           \right)[0{:}8]\;\text{as } u64$
    \State $s_{\mathrm{rev}} \gets \mathrm{SHA256}\!\left(\texttt{"SBF\_REV\_"} \,\|\,
           t^{\mathrm{le32}} \,\|\, \texttt{"\_"} \,\|\, s_i^{\mathrm{le64}}
           \right)[0{:}8]\;\text{as } u64$
    \If{$\mathrm{chain.step}(\mathrm{rng\_fwd}(s_{\mathrm{fwd}}),\,
         \mathrm{rng\_rev}(s_{\mathrm{rev}}))$ accepted}
      \State $\mathrm{accepted}_i \gets \mathrm{accepted}_i + 1$
    \EndIf
  \EndFor
  \State $\mathrm{endpoints}.\mathrm{push}\!\left(
         \ec(\mathrm{chain.assignment}),\; i,\; \mathrm{chain.assignment}
         \right)$
  \State $\pi_{\mathrm{cur}} \gets \mathrm{chain.assignment}$
         \Comment{restart next burst from current endpoint}
  \State \textbf{Record} $\alpha_i = \mathrm{accepted}_i / \ell$ in audit manifest
\EndFor
\State Sort $\mathrm{endpoints}$ by $\ec$ ascending, then by $i$ ascending
\State \Return $\mathrm{endpoints}[\lfloor p \cdot n \rfloor].\mathrm{assignment}$
\end{algorithmic}
\end{algorithm}

\paragraph{Two RNGs per step.}
Each ForestRecom step requires two independent RNG streams:
the \emph{forward RNG} (to sample the proposed spanning tree and cut)
and the \emph{reverse RNG} (to evaluate the reverse proposal probability for
the Metropolis--Hastings ratio).
Both are derived deterministically from the burst seed $s_i$ and the
step index $t$ using the \texttt{SBF\_FWD\_} and \texttt{SBF\_REV\_}
prefixes.
This ensures that any burst trajectory can be independently reproduced by
any verifier who holds the base seed $b$ and the burst parameters.

\paragraph{Endpoint restart.}
At the end of each burst, $\pi_{\mathrm{cur}}$ is updated to
\texttt{chain.assignment}, the endpoint of the current burst.
This is unconditional: the burst endpoint always becomes the next starting
plan, regardless of whether it has lower or higher edge-cut than the
previous $\pi_{\mathrm{cur}}$.
This allows the chain to escape local minima by occasionally accepting
worse endpoints during the burst (through the MH criterion).

\paragraph{CLI invocation.}
\begin{verbatim}
redist state --state NC --year 2020 \
    --search short-burst-forest \
    --burst-length 20 --n-bursts 50 --percentile 0.0
\end{verbatim}

\subsection{ShortBurstMergeSplit}
\label{subsec:alg-sbms}

\sbms{} is identical to Algorithm~\ref{alg:sbf} except:
\begin{enumerate}
  \item The inner chain is MergeSplitChain instead of ForestRecomChain.
  \item The seed prefixes are \texttt{"SBMS\_CHAIN\_"}, \texttt{"SBMS\_FWD\_"},
        and \texttt{"SBMS\_REV\_"}.
  \item The MH ratio is a single spanning-tree-count ratio rather than the
        two-tree ForestRecom ratio, so the reverse RNG is used differently
        internally (see G.10), but the interface to the burst loop is
        identical.
\end{enumerate}

\paragraph{CLI invocation.}
\begin{verbatim}
redist state --state NC --year 2020 \
    --search short-burst-merge-split \
    --burst-length 20 --n-bursts 50 --percentile 0.0
\end{verbatim}

\subsection{Per-Burst Acceptance Rate Diagnostic}
\label{subsec:alg-acceptance}

\begin{definition}[Per-burst acceptance rate]
\label{def:alpha-burst}
For burst $i$ with $\ell$ steps of which $a_i$ are accepted, the per-burst
acceptance rate is
\[
  \alpha_i = \frac{a_i}{\ell} \in [0,1].
\]
The mean per-burst acceptance rate over $n$ bursts is
$\bar\alpha = \frac{1}{n}\sum_{i=0}^{n-1} \alpha_i$.
\end{definition}

$\alpha_i$ is recorded in the run audit manifest alongside $\ec$ of the
burst endpoint and the burst seed $s_i$.
It serves two diagnostic purposes.

\paragraph{Effective burst length.}
Only accepted steps produce a new plan; rejected steps stay at the current
plan.
The effective number of distinct plans visited per burst is approximately
$\alpha_i \cdot \ell$.
For standard \sburst{} with $\bar\alpha \approx 0.75$, a burst of length 20
visits roughly 15 distinct plans.
For \sbf{} with $\bar\alpha \approx 0.45$, the same 20-step burst visits
roughly 9 distinct plans.
Lower effective burst length directly limits the burst's ability to find
a substantially better plan than the starting $\pi_{\mathrm{cur}}$.

\paragraph{Anomaly detection.}
Unusually low $\alpha_i$ for a specific burst (e.g., $\alpha_i < 0.1$)
may indicate that the starting plan for that burst was in a region of the
plan space with very few valid neighbouring proposals.
Such anomalous bursts contribute little to the endpoint distribution and
should be investigated if they account for a large fraction of the total
compute budget.

\begin{remark}[Acceptance rate is not a quality metric]
\label{rem:accept-not-quality}
Higher $\bar\alpha$ is \emph{not} intrinsically better.
Standard \sburst{} has high $\bar\alpha$ because it makes no attempt at
distributional correctness; \sbf{} has lower $\bar\alpha$ because it is
doing additional computational work (evaluating the MH ratio) to ensure
that accepted plans are approximately correctly weighted.
The tradeoff is between exploration speed and distributional fidelity,
not between ``good'' and ``bad'' acceptance rates.
\end{remark}
